\documentclass[aspectratio=169,11pt]{beamer}
\usepackage{paperdeck}
\title[ALiBi]{Train Short, Test Long: Attention with Linear Biases Enables Input Length Extrapolation}
\subtitle{Press, Smith and Lewis (ICLR 2022), arXiv:2108.12409}
\author[Paper summary]{Paper summary for the CS-781 capstone}
\date{}
\begin{document}

\begin{frame}\titlepage\end{frame}

\paperglance
{O.\ Press, N.\ A.\ Smith, M.\ Lewis. \emph{Train Short, Test Long: Attention with Linear Biases Enables Input Length Extrapolation}. ICLR 2022. University of Washington, Facebook AI Research, Allen Institute for AI.}
{Drop position embeddings entirely. Subtract a fixed, head-specific penalty proportional to query--key \textbf{distance} from each attention score. Models trained on short inputs then work on longer ones.}
{Our A3 prefix-LM uses ALiBi and has no position embeddings. Legal judgments are long, so extrapolation matters for Phases 4--7.}
{Main deck: ``A3: hybrid prefix-LM with ALiBi''. Glossary: ``Positional information (3)''.}

\begin{frame}{Roadmap}
  \begin{multicols}{2}\small\tableofcontents\end{multicols}
\end{frame}

% ------------------------------------------------------------------
\section{The problem: extrapolation}

\begin{frame}{Can a model read longer inputs than it was trained on?}
  \begin{columns}[T]
    \begin{column}{0.52\textwidth}
      \small
      \textbf{Extrapolation}: performing well when the input at inference is longer than the training length $L$.
      \vspace{1mm}

      Why we want it:
      \begin{itemize}
        \item Attention costs $O(L^2)$, so training on long sequences is expensive.
        \item If a model could train on $L=512$ and still read 2,048 tokens later, training would be much cheaper.
      \end{itemize}
      Vaswani et al.\ speculated that sinusoidal positions ``may'' extrapolate. \textbf{Nobody had checked carefully.}
    \end{column}
    \begin{column}{0.44\textwidth}
      \begin{exampleblock}{Legal motivation}
        \small
        Our Phase-1 models train with a 1,024-token window. A typical Supreme Court judgment is many thousands of tokens.
        For judgment prediction or case retrieval, we would like to feed more of the document at inference without retraining on long sequences.
      \end{exampleblock}
    \end{column}
  \end{columns}
\end{frame}

\begin{frame}{Existing position methods, tested for extrapolation}
  \begin{columns}[T]
    \begin{column}{0.54\textwidth}
      \small
      Setup: WikiText-103 language model (Baevski and Auli, 2018: 16 layers, $d=1024$, 8 heads). Train with $L=512$ or $1024$; evaluate on longer inputs, changing \emph{only} the position method.
      \vspace{1mm}

      \begin{itemize}
        \item \textbf{Learned positions}: no embedding exists beyond $L$, so extrapolation is impossible.
        \item \textbf{Sinusoidal}: improves for up to about 20--50 extra tokens, then degrades.
        \item \textbf{Rotary (RoPE)}: better; gains up to about 100--200 extra tokens. It was slower in their implementation.
        \item \textbf{T5 bias}: best; gains up to about 600--800 extra tokens, but training is at least \textbf{twice as slow}.
      \end{itemize}
    \end{column}
    \begin{column}{0.42\textwidth}
      \begin{exampleblock}{Key finding}
        \small
        Extrapolation depends heavily on the \textbf{position method}.
        \vspace{1mm}

        Methods that inject position \emph{at every layer into attention}, and not into the values, extrapolate better. ALiBi builds on this observation.
      \end{exampleblock}
      \begin{alertblock}{No efficiency gain from T5 bias}
        \small
        T5 bias at $L=512$ runs at 14.4k words/s; sinusoidal at $L=1024$ runs at 28.5k. Training short with T5 bias saves nothing.
      \end{alertblock}
    \end{column}
  \end{columns}
\end{frame}

% ------------------------------------------------------------------
\section{The basic idea: linear biases}

\begin{frame}{ALiBi: attention with linear biases}
  \begin{columns}[T]
    \begin{column}{0.52\textwidth}
      \small
      No position embeddings anywhere. For query $i$ attending to keys $1\ldots i$, add a fixed bias:
      \[
        \text{softmax}\big(q_iK^\top + m\cdot[-(i-1),\ldots,-2,-1,0]\big)
      \]
      \begin{itemize}
        \item Penalty grows \textbf{linearly} with distance.
        \item $m$ is a \textbf{head-specific slope}, fixed, not learned.
        \item Position enters every layer, through attention scores only; values carry no position.
      \end{itemize}
      \vspace{1mm}
      \textbf{Implementation}: add the bias to the attention mask, a few lines of code. There is no runtime penalty; the only extra is a per-head mask ($n\times L\times L$, up to about 100 MB).
    \end{column}
    \begin{column}{0.44\textwidth}
      \centering
      \small
      Bias matrix for one head, $m=1$:
      \vspace{1mm}

      \begin{tabular}{@{}c|ccccc@{}}
        & $k_1$ & $k_2$ & $k_3$ & $k_4$ & $k_5$\\
        \hline
        $q_1$ & 0 & & & &\\
        $q_2$ & $-1$ & 0 & & &\\
        $q_3$ & $-2$ & $-1$ & 0 & &\\
        $q_4$ & $-3$ & $-2$ & $-1$ & 0 &\\
        $q_5$ & $-4$ & $-3$ & $-2$ & $-1$ & 0\\
      \end{tabular}
      \vspace{1mm}

      {\scriptsize Multiply by the head's slope $m$ and add to $q_i\cdot k_j$. Blank cells are masked (future).}
    \end{column}
  \end{columns}
\end{frame}

\begin{frame}{Choosing the slopes}
  \begin{columns}[T]
    \begin{column}{0.52\textwidth}
      \small
      For $n$ heads, the slopes form a geometric sequence that starts at $2^{-8/n}$ with ratio $2^{-8/n}$.
      \begin{itemize}
        \item 8 heads: $\tfrac12,\tfrac14,\tfrac18,\ldots,\tfrac1{256}$.
        \item 16 heads: $\tfrac{1}{2^{0.5}},\tfrac12,\tfrac{1}{2^{1.5}},\ldots,\tfrac{1}{2^8}$.
      \end{itemize}
      \vspace{1mm}
      Findings from about ten slope sets tried:
      \begin{itemize}
        \item The best sets have slopes in $(0,1)$, denser near 0.
        \item The method is robust to the exact choice; the same slopes worked across domains and sizes.
        \item \textbf{Trainable slopes} extrapolated poorly and slowed training by 3\%.
      \end{itemize}
    \end{column}
    \begin{column}{0.44\textwidth}
      \begin{exampleblock}{Different heads, different ranges}
        \small
        \begin{itemize}
          \item Slope $\tfrac12$: a key 10 tokens back loses 5 points. A \textbf{local} head, for grammar.
          \item Slope $\tfrac{1}{256}$: 10 tokens back loses only 0.04. A \textbf{long-range} head, for topic, parties and sections.
        \end{itemize}
      \end{exampleblock}
      \begin{block}{Our A3 (12 heads)}
        \small
        12 is not a power of two. As the ALiBi code does, we take the 8-head slopes plus every other 16-head slope: 0.707, 0.354, 0.177 and 0.088.
      \end{block}
    \end{column}
  \end{columns}
\end{frame}

\begin{frame}{Easy worked example}
  \begin{columns}[T]
    \begin{column}{0.52\textwidth}
      \small
      Query at position 5; keys at positions 1--5 with raw scores $q\cdot k=[2.0,\;1.0,\;1.5,\;0.5,\;1.0]$. Distances $[4,3,2,1,0]$.
      \vspace{1mm}

      \footnotesize
      \begin{tabular}{@{}lccccc@{}}
        \toprule
        Key position & 1 & 2 & 3 & 4 & 5\\
        \midrule
        Raw weight (no bias) & \textbf{0.39} & 0.14 & 0.24 & 0.09 & 0.14\\
        Bias, $m=0.5$ & $-2$ & $-1.5$ & $-1$ & $-0.5$ & 0\\
        Biased score & 0.0 & $-0.5$ & 0.5 & 0.0 & 1.0\\
        ALiBi weight & 0.14 & 0.09 & 0.24 & 0.14 & \textbf{0.39}\\
        \bottomrule
      \end{tabular}
    \end{column}
    \begin{column}{0.44\textwidth}
      \begin{exampleblock}{What happened}
        \small
        With slope 0.5, the farthest key drops from 0.39 to 0.14, and the nearest key rises from 0.14 to 0.39.
        A strongly matching distant key can still win, but it must match \emph{much} better.
        \vspace{1mm}

        With slope $1/256$ the weights barely change: that head still looks far back.
      \end{exampleblock}
      \begin{block}{Why it extrapolates}
        \small
        A distance of 1,500 needs no new parameters: the penalty is just $1500\,m$. Very distant keys are simply faded out.
      \end{block}
    \end{column}
  \end{columns}
\end{frame}

% ------------------------------------------------------------------
\section{Results}

\begin{frame}{WikiText-103: better, even without extrapolating}
  \begin{columns}[T]
    \begin{column}{0.52\textwidth}
      \small
      \begin{itemize}
        \item For every training length $L$, ALiBi beats sinusoidal \emph{at the same length}. For example, at $L=3072$: \textbf{17.60} vs 18.67 perplexity.
        \item ALiBi trained on $L=512$ scores 19.73 at 512 tokens, and \textbf{18.40} when evaluated on 3,072 tokens. That beats the sinusoidal model \emph{trained} on 3,072.
        \item That $L=512$ model trains \textbf{1.84$\times$ faster} and fits on a smaller GPU.
        \item At $L=1024$ and $3072$, ALiBi also beats rotary and T5 bias without extrapolation.
        \item The $L=512$ model keeps improving up to inputs of over 12k tokens.
      \end{itemize}
    \end{column}
    \begin{column}{0.44\textwidth}
      \begin{exampleblock}{Books domain}
        \small
        The same slopes, unchanged, were used on the Toronto Book Corpus. Results transferred fully: better than sinusoidal both with and without extrapolation.
      \end{exampleblock}
      \begin{block}{Recency bias}
        \small
        ALiBi prefers nearby tokens. For language modelling on small data this inductive bias helps even at equal length.
      \end{block}
    \end{column}
  \end{columns}
\end{frame}

\begin{frame}{Scaling up: 1.3B parameters on 461 GB}
  \begin{columns}[T]
    \begin{column}{0.52\textwidth}
      \small
      Model: 25 layers, 16 heads, $d=2048$, FFN 8192: \textbf{1.3B parameters}. Data: RoBERTa corpus + English CC-100 (461 GB). One epoch = 50k updates on 128 V100 GPUs.
      \vspace{1mm}

      \begin{itemize}
        \item ALiBi trained on $L=1024$ \textbf{beats} sinusoidal trained on $L=2048$ at 2,048 tokens by 0.09 perplexity, using 3.1 GB less memory.
        \item It reaches a given perplexity about \textbf{11\% faster}.
        \item ALiBi $L=512$ vs sinusoidal $L=1024$ at 1,024 tokens: only 0.06 behind, 7\% faster, 1.6 GB less memory.
        \item At equal $L$, ALiBi and sinusoidal are \emph{similar} at this scale.
      \end{itemize}
    \end{column}
    \begin{column}{0.44\textwidth}
      \begin{exampleblock}{The ``2$L$'' sweet spot}
        \small
        On this corpus, the best perplexity comes at about \textbf{twice} the training length:
        \begin{itemize}
          \item $L=512$: 9.79 at 512, best 9.3 at 1,012 tokens.
          \item $L=1024$: 9.16 at 1,024, best 8.9 at 2,024 tokens.
        \end{itemize}
        Performance stays strong up to 10,000 tokens. Sinusoidal degrades as soon as one token is added beyond $L$.
      \end{exampleblock}
    \end{column}
  \end{columns}
\end{frame}

\begin{frame}{Why does longer input help? The early-token curse}
  \begin{columns}[T]
    \begin{column}{0.52\textwidth}
      \small
      Evaluation usually splits text into non-overlapping chunks of length $L_{\text{valid}}$. Tokens at the \emph{start} of each chunk have little context, so they get worse predictions: the \textbf{early-token curse}.
      \vspace{1mm}

      Two possible reasons ALiBi improves with longer inputs:
      \begin{enumerate}
        \item it really uses context more than $L$ tokens back; or
        \item longer chunks mean fewer tokens suffer the early-token curse.
      \end{enumerate}
      Using sliding-window evaluation, the authors find that ALiBi's gain is \textbf{largely explained by reason 2}.
    \end{column}
    \begin{column}{0.44\textwidth}
      \begin{alertblock}{An honest caveat}
        \small
        ALiBi lets a model \emph{accept} long inputs gracefully. It does not prove the model \emph{uses} information from thousands of tokens back.
        The authors suggest future work should exploit longer histories more efficiently.
      \end{alertblock}
      \begin{exampleblock}{For our evaluation}
        \small
        Our continuation loss always gives the model 768 tokens of context, so it avoids the early-token curse by design.
      \end{exampleblock}
    \end{column}
  \end{columns}
\end{frame}

% ------------------------------------------------------------------
\section{Assessment and outlook}

\begin{frame}{ALiBi in our capstone}
  \begin{columns}[T]
    \begin{column}{0.5\textwidth}
      \small
      \begin{itemize}
        \item A3 uses ALiBi with 12 heads and no position embeddings.
        \item \textbf{Our extension}: the paper uses causal attention only. A3 also has a \emph{bidirectional} prefix, so we use the symmetric penalty $-m\,|i-j|$ inside the prompt.
        \item The bias is built as a float mask together with the prefix mask. Our tests check it against cached generation.
      \end{itemize}
    \end{column}
    \begin{column}{0.46\textwidth}
      \begin{block}{Experiments worth running}
        \small
        \begin{itemize}
          \item Evaluate A1 (learned), A3 (ALiBi) and A4 (RoPE) at 1,024, 1,536 and 2,048 tokens. A1 cannot go beyond 1,024 at all.
          \item Train A3 at $L=512$ and test at 1,024: can it match a 1,024-trained model, as in the paper?
          \item For long judgments in Phases 4--7, compare ALiBi extrapolation against chunking.
        \end{itemize}
      \end{block}
    \end{column}
  \end{columns}
\end{frame}

\begin{frame}{Strengths and limitations}
  \begin{columns}[T]
    \begin{column}{0.48\textwidth}
      \begin{exampleblock}{Strengths}
        \small
        \begin{itemize}
          \item Extremely simple; no parameters; no runtime cost.
          \item Real extrapolation: train short, test long.
          \item Saves training time and memory at a given test length.
          \item Better perplexity than sinusoidal on smaller data.
          \item Slopes transfer across domains and model sizes.
          \item Careful, controlled comparison of four position methods.
        \end{itemize}
      \end{exampleblock}
    \end{column}
    \begin{column}{0.48\textwidth}
      \begin{alertblock}{Limitations}
        \small
        \begin{itemize}
          \item Best results come at about 2$L$; much longer inputs help less.
          \item Gains largely come from avoiding the early-token curse, not from long-range use.
          \item The recency bias may hurt tasks that need precise distant retrieval.
          \item Designed and tested for causal language modelling only.
          \item At billion-parameter scale, no gain at equal length.
          \item Slopes are hand-chosen.
        \end{itemize}
      \end{alertblock}
    \end{column}
  \end{columns}
\end{frame}

\begin{frame}{Future research directions}
  \begin{columns}[T]
    \begin{column}{0.5\textwidth}
      \small
      \textbf{From the paper}
      \begin{itemize}
        \item Exploit longer histories more efficiently, beyond reducing the early-token curse.
        \item Combine ALiBi with orthogonal methods such as kNN-LM or Routing Transformer.
        \item Use the saved memory to stack more layers.
      \end{itemize}
      \textbf{What followed}
      \begin{itemize}
        \item Adopted in large open models such as BLOOM (BigScience, 2022) and MPT.
        \item Length-extrapolation research: RoPE scaling methods (Position Interpolation, YaRN), xPos, and studies of positional schemes such as NoPE (Kazemnejad et al., 2023).
      \end{itemize}
    \end{column}
    \begin{column}{0.46\textwidth}
      \begin{block}{Open questions}
        \small
        \begin{itemize}
          \item Can a model \emph{use} information thousands of tokens back, not just tolerate it?
          \item Should slopes adapt to the task or the domain?
          \item How does recency bias interact with retrieval of a single far-away fact (``needle in a haystack'')?
          \item For legal text: parties and statutes named on page 1 matter on page 40. Is a linear recency penalty the right prior for judgments?
        \end{itemize}
      \end{block}
    \end{column}
  \end{columns}
\end{frame}

\begin{frame}{Summary}
  \begin{itemize}
    \item Sinusoidal positions do not extrapolate; rotary helps a little; T5 bias helps more but is slow.
    \item \textbf{ALiBi}: no position embeddings; add $-m\cdot$distance to each attention score, with fixed geometric slopes per head.
    \item Worked example: slope 0.5 moves attention from the farthest key (0.39) to the nearest (0.39).
    \item Results: better perplexity at equal length on WikiText-103; a model trained on 512 tokens beats a sinusoidal model trained on 3,072; at 1.3B, train on 1,024 and match or beat 2,048, 11\% faster.
    \item Best around 2$L$; gains mostly from reducing the early-token curse.
    \item For us: the position scheme of A3, extended to a bidirectional prefix.
  \end{itemize}
\end{frame}

\begin{frame}{Reference}
  \small
  Press, O., Smith, N.\ A., and Lewis, M.\ (2022). Train Short, Test Long: Attention with Linear Biases Enables Input Length Extrapolation. In \emph{International Conference on Learning Representations (ICLR)}. arXiv:2108.12409.
  \vspace{3mm}

  {\footnotesize Other work cited: Vaswani et al.\ (2017); Raffel et al.\ (2020) T5; Su et al.\ (2021) RoFormer; Press et al.\ (2021) Shortformer; BigScience (2022) BLOOM; Chen et al.\ (2023) Position Interpolation; Peng et al.\ (2023) YaRN; Sun et al.\ (2023) xPos; Kazemnejad et al.\ (2023) The Impact of Positional Encoding on Length Generalization in Transformers.}
\end{frame}

\end{document}
